% 图 1.2 GAN-based JSCC 结构图（重绘加强版）
% 基于 CICC0903540 技术文档 2.2 - 语义通信技术原理
\documentclass[tikz,border=10pt]{standalone}

% 强制全局样式与配色
\usepackage{amssymb, amsmath}
\usepackage{fontspec}
\usepackage{xeCJK}
\setCJKmainfont{SimHei}
\usepackage{tikz}
\usetikzlibrary{
  arrows.meta, positioning, shapes.geometric, calc,
  backgrounds, fit, decorations.pathreplacing, patterns
}

% 低饱和度马卡龙配色（禁止使用纯色）
\definecolor{mBlue}{HTML}{AEC6CF}
\definecolor{mPink}{HTML}{FFD1DC}
\definecolor{mGreen}{HTML}{B1D8B7}
\definecolor{mPurple}{HTML}{B39EB5}
\definecolor{mYellow}{HTML}{FDFD96}
\definecolor{mGray}{HTML}{E5E4E2}
\definecolor{mDark}{HTML}{4A4A4A}

% 全局 TikZ 样式
\tikzset{
  >=stealth,
  thick,
  draw=mDark,
  % 基础节点样式
  block/.style={
    rectangle, rounded corners=3pt, draw=mDark, align=center,
    font=\sffamily\small, minimum height=0.9cm, inner sep=5pt},
  % 微观模块样式
  micro/.style={
    rectangle, rounded corners=2pt, draw=mDark!70, align=center,
    font=\sffamily\tiny, minimum height=0.4cm, inner sep=3pt},
  % 数学操作符节点
  op/.style={
    circle, draw=mDark, minimum size=0.45cm, inner sep=0,
    font=\scriptsize, fill=mGray!60},
  % 张量维度标注
  ts/.style={font=\tiny\sffamily, color=mDark!80},
  % 箭头样式
  arr/.style={-{Stealth[length=4pt]}, thick, draw=mDark},
  darr/.style={-{Stealth[length=4pt]}, thick, draw=mDark, dashed},
  % 图例
  legendbox/.style={
    rectangle, draw=mDark!50, rounded corners=3pt,
    fill=mGray!15, inner sep=5pt, font=\tiny\sffamily},
}

\begin{document}
\begin{tikzpicture}[
  node distance=0.7cm and 1.2cm,
]
  % ========== 主数据流（发送端 → 信道 → 接收端）==========
  % 输入图像
  \node[block, fill=mGray!50, minimum width=1.5cm] (img) {$\mathbf{x}$\\{\tiny 输入}};

  % Encoder（带微观结构拆解）
  \node[block, fill=mBlue!40, right=1.2cm of img, minimum width=2cm] (enc)
    {Encoder\\{\tiny 语义特征提取}};

  % 功率控制
  \node[block, fill=mBlue!40, right=1.2cm of enc, minimum width=1.8cm] (pwr)
    {功率控制\\{\tiny Power Ctrl}};

  % AWGN 信道
  \node[block, fill=mGreen!50, right=1.2cm of pwr, minimum width=1.8cm] (awgn)
    {AWGN\\{\tiny 加性高斯白噪声}};

  % Generator（带微观结构拆解）
  \node[block, fill=mPurple!40, right=1.2cm of awgn, minimum width=2cm] (gen)
    {Generator\\{\tiny 图像重建}};

  % 输出图像
  \node[block, fill=mGray!50, right=1.2cm of gen, minimum width=1.5cm] (out)
    {$\hat{\mathbf{x}}$\\{\tiny 重建}};

  % ========== 主流程连线 + 张量维度标注 ==========
  \draw[arr] (img) -- (enc)
    node[pos=0.5, above=3pt, ts] {$1{\times}3{\times}256{\times}256$}
    node[pos=0.5, below=3pt, ts] {RGB 原始图像};

  \draw[arr] (enc) -- (pwr)
    node[pos=0.5, above=3pt, ts] {$\mathbf{y} \in \mathbb{R}^{1{\times}32{\times}32{\times}32}$}
    node[pos=0.5, below=3pt, ts] {语义特征向量};

  \draw[arr] (pwr) -- (awgn)
    node[pos=0.5, above=3pt, ts] {$\sqrt{P}\mathbf{y}$};

  \draw[arr] (awgn) -- (gen)
    node[pos=0.5, above=3pt, ts] {$\hat{\mathbf{y}} = \sqrt{P}\mathbf{y} + \mathbf{n}$}
    node[pos=0.5, below=3pt, ts] {受噪声污染的特征};

  \draw[arr] (gen) -- (out)
    node[pos=0.5, above=3pt, ts] {$1{\times}3{\times}256{\times}256$}
    node[pos=0.5, below=3pt, ts] {重建图像};

  % ========== Discriminator（仅训练阶段使用，用虚线）==========
  \node[block, fill=mPink!40, below=1.8cm of awgn, dashed, minimum width=2.2cm] (disc)
    {Discriminator\\{\tiny 判别器（仅训练）}};

  % 从真实图像 x 到判别器
  \draw[darr] (img.south) -- ++(0,-0.4) -| (disc.west)
    node[pos=0.25, left=2pt, ts] {$\mathbf{x}$};

  % 从重建图像 x_hat 到判别器
  \draw[darr] (out.south) -- ++(0,-0.4) -| (disc.east)
    node[pos=0.25, right=2pt, ts] {$\hat{\mathbf{x}}$};

  % 从信道输出 y_hat 到判别器（条件输入）
  \draw[darr] (awgn.south) -- (disc.north)
    node[pos=0.5, right=2pt, ts] {$\hat{\mathbf{y}}$\\{\tiny 条件}};

  % 判别器输出
  \node[below=0.4cm of disc, ts] {$D(\mathbf{x}, \hat{\mathbf{y}}) \in [0,1]$};

  % ========== Encoder 微观结构拆解 ==========
  \node[micro, fill=mBlue!30, below=1.0cm of enc, xshift=-1.2cm] (e1) {Conv};
  \node[micro, fill=mBlue!30, right=0.15cm of e1] (e2) {BatchNorm};
  \node[micro, fill=mBlue!30, right=0.15cm of e2] (e3) {ReLU};
  \node[micro, fill=mBlue!25, right=0.2cm of e3, minimum width=1.2cm] (e4) {ResBlock$\times$2};

  \draw[arr, -, thick] (e1) -- (e2) -- (e3) -- (e4);
  \draw[arr, dashed, mDark!60] (enc.south) |- (e2.north);

  \node[ts, below=0.15cm of e1, anchor=north] {$3{\to}64$};
  \node[ts, below=0.15cm of e4, anchor=north] {$256{\to}32$};

  % ========== Generator 微观结构拆解 ==========
  \node[micro, fill=mPurple!30, below=1.0cm of gen, xshift=-1.2cm] (g1) {Conv};
  \node[micro, fill=mPurple!30, right=0.15cm of g1] (g2) {BatchNorm};
  \node[micro, fill=mPurple!30, right=0.15cm of g2] (g3) {ReLU};
  \node[micro, fill=mPurple!25, right=0.2cm of g3, minimum width=1.2cm] (g4) {ResBlock$\times$5};

  \draw[arr, -, thick] (g1) -- (g2) -- (g3) -- (g4);
  \draw[arr, dashed, mDark!60] (gen.south) |- (g2.north);

  \node[ts, below=0.15cm of g1, anchor=north] {$32{\to}256$};
  \node[ts, below=0.15cm of g4, anchor=north] {$256{\to}3$};

  % ========== 层次化表达：背景分组 ==========
  \begin{scope}[on background layer]
    % 发送端
    \node[fit=(img)(enc)(pwr)(e1)(e4),
      fill=mBlue!12, draw=mDark!40, dashed,
      rounded corners=5pt, inner sep=8pt] {};
    \node[above=0.3cm of enc, font=\sffamily\footnotesize\bfseries, color=mDark]
      {发送端（编码器）};

    % 接收端
    \node[fit=(gen)(out)(g1)(g4),
      fill=mPurple!12, draw=mDark!40, dashed,
      rounded corners=5pt, inner sep=8pt] {};
    \node[above=0.3cm of gen, font=\sffamily\footnotesize\bfseries, color=mDark]
      {接收端（解码器）};
  \end{scope}

  % ========== 信息密度拉满：损失函数标注 ==========
  \node[legendbox, anchor=north west] at ([xshift=5pt, yshift=-5pt]current bounding box.north west)
    [align=left, text width=4.5cm] {
    \textbf{训练损失函数}\\[2pt]
    $L_t = \lambda L_G + \alpha L_{\text{MSE}} + \beta L_{\text{LPIPS}}$\\[1pt]
    $L_G = -\mathbb{E}[\log D(G(\hat{\mathbf{y}}), \hat{\mathbf{y}})]$\\[1pt]
    $L_D = -\mathbb{E}[\log D(\mathbf{x}, \mathbf{y})] - \mathbb{E}[\log(1 - D(\hat{\mathbf{x}}, \hat{\mathbf{y}}))]$
  };

  % ========== 图例（右下角）==========
  \node[legendbox, anchor=south east] at ([xshift=-5pt, yshift=5pt]current bounding box.south east)
    [align=left, text width=4.2cm] {
    \textbf{图例}\\[1pt]
    \colorbox{mBlue!40}{\rule{0.16cm}{0.12cm}\hspace{0.05cm}} 编码端 \quad
    \colorbox{mPurple!40}{\rule{0.16cm}{0.12cm}\hspace{0.05cm}} 解码端 \quad
    \colorbox{mGreen!50}{\rule{0.16cm}{0.12cm}\hspace{0.05cm}} 信道\\[1pt]
    实线：推理路径 \quad 虚线：训练路径\\[1pt]
    SNR = $10\log_{10}\frac{\|\mathbf{y}\|^2}{\sigma^2}$ (dB)
  };
\end{tikzpicture}
\end{document}
